% ===========================================================
\section{Methods}
\label{sec:methods}
% ===========================================================

\subsection{Graph Representation}

Each crystal structure is converted to a directed graph
$\mathcal{G}=(\mathcal{V},\mathcal{E})$ where nodes $v_i\in\mathcal{V}$
represent atoms and edges $(i,j)\in\mathcal{E}$ connect pairs
within a cutoff radius $r_c = 6$\,\AA.
The node feature $\mathbf{h}_i^{(0)}\in\mathbb{R}^{100}$ is the
per-atom chemical descriptor.
Each edge carries a geometric feature comprising
(i) a radial basis function (RBF) expansion of the interatomic
distance $d_{ij}$
\begin{equation}
  e_k(d_{ij}) = \exp\!\left(-\gamma\bigl(d_{ij}-\mu_k\bigr)^2\right),
  \quad k=1,\ldots,64,
\end{equation}
and (ii) the unit direction vector $\hat{\mathbf{r}}_{ij}$.

\subsection{Equivariant Message Passing}

Node features are updated through $L=5$ message-passing layers
\begin{equation}
  \mathbf{m}_{ij} = \sigma\!\left(
    W_m[\mathbf{h}_i\,\|\,\mathbf{h}_j\,\|\,\mathbf{e}(d_{ij})\,\|
    \langle\mathbf{h}_i,\hat{\mathbf{r}}_{ij}\rangle]
  \right) \odot g(\mathbf{h}_i,\mathbf{e}),
\end{equation}
\begin{equation}
  \mathbf{h}_i^{(l+1)} = \mathbf{h}_i^{(l)}
    + W_u\!\left[\mathbf{h}_i^{(l)}\,\Bigl\|
      \sum_{j\in\mathcal{N}(i)}\mathbf{m}_{ij}\right],
\end{equation}
where $g$ is a learned scalar gate, $\sigma$ is the SiLU activation,
and residual connections are used throughout.

\subsection{Graph Pooling and Output Heads}

A graph-level embedding is obtained via attention pooling.
Two task-specific heads are then applied:
(1) a \emph{phonon head} producing the flattened band structure
    $\hat{\omega}\in\mathbb{R}^{N_q\times 24}$, and
(2) an \emph{elastic head} producing
    $\hat{\mathbf{C}}\in\mathbb{R}^{12}$.

\subsection{Physics-Informed Training Objective}
\label{sec:loss}

The total loss is
\begin{equation}
  \mathcal{L} = \mathcal{L}_{\rm ph}
              + \lambda_s\,\mathcal{L}_{\rm smooth}
              + \lambda_e\,\mathcal{L}_{\rm elastic},
  \label{eq:loss}
\end{equation}
where
\begin{align}
  \mathcal{L}_{\rm ph}     &= \frac{1}{N_q B}\sum_{q,b}
       (\omega_{qb}-\hat\omega_{qb})^2, \\
  \mathcal{L}_{\rm smooth} &= \frac{1}{N_q B}\sum_{q,b}
       (\hat\omega_{q+1,b}-2\hat\omega_{q,b}+\hat\omega_{q-1,b})^2,\\
  \mathcal{L}_{\rm elastic}&= \frac{1}{12}\sum_k(C_k-\hat C_k)^2.
\end{align}
We use $\lambda_s=1.5$ and $\lambda_e=0.2$.

\subsection{Training Details}

The model is trained with AdamW ($\eta=3\times10^{-4}$,
weight decay $10^{-5}$) for up to 600 epochs with cosine
learning-rate annealing and early stopping (patience = 100).
All experiments are run on a single Kaggle GPU (P100).

\subsection{Feature Importance Ablation}
\label{sec:ablation}

To quantify the contribution of each atomic descriptor we perform
a leave-one-out ablation study using XGBoost as a fast surrogate.
For each feature $f$ we retrain the model with $f$ removed and
record $\Delta{\rm MAE}=\text{MAE}_{-f}-\text{MAE}_{\rm full}$;
positive values indicate the feature is beneficial.
